\section{结论与展望}

\subsection{工作总结}
本文围绕“动脉瘤精确检测与破裂风险预测”任务，完成了从算法实现到工程流程贯通的毕业设计工作。主要成果如下：
\begin{enumerate}
	\item 基于三维Attention U-Net完成了颅内动脉瘤分割模块，建立了病例级流式训练、补丁采样和整例重建推理流程；
	\item 采用Dice与BCE联合损失提升前景学习稳定性，并实现了概率图/二值图双输出保存机制；
	\item 针对风险预测任务，构建了临床与组学特征预处理流水线，完成清洗、筛选、编码与不平衡处理；
	\item 设计并实现了双流交叉注意力风险模型，实现破裂风险二分类训练、阈值优化和测试评估；
	\item 形成了统一配置驱动的训练脚本体系，满足了任务书中“精确检测+风险预测+实验验证”的要求。
\end{enumerate}

综合来看，本文实现方案具备较好的工程可复现性和模块化扩展能力，能够为后续临床辅助决策系统开发提供基础。

\subsection{不足与后续展望}
尽管本文完成了既定目标，但仍有进一步改进空间：
\begin{itemize}
	\item 数据规模与多中心验证不足，后续应在更大样本和跨设备数据上验证泛化性能；
	\item 分割模型可进一步引入Transformer或自监督预训练策略，提高复杂血管背景下鲁棒性；
	\item 风险预测可结合分割网络深层特征，实现端到端多模态联合优化；
	\item 临床可解释性仍需增强，后续可增加特征归因分析与病例级可视化报告；
	\item 系统部署方面可进一步完善接口标准与推理加速策略，提升真实场景可用性。
\end{itemize}

总体而言，本文工作验证了深度学习方法在颅内动脉瘤自动化分析中的应用潜力，也为后续研究与工程落地提供了清晰起点。
